\documentclass{report}
\usepackage{amsfonts, amsmath}
\newcommand\R{{\mathbb R}}
\newcommand\Q{{\mathbb Q}}
\newcommand\N{{\mathbb N}}
\newcommand\Or{{\mathcal O}}
\newcommand\ve{\varepsilon}
%\pagestyle{empty}
\begin{document}
\large
\centerline{\Huge Analisi Matematica II}
\renewcommand{\thefootnote}{ } 
\renewcommand{\thefootnote}{\arabic{footnote}} 
\vskip.1cm
\centerline{\Large Soluzioni esonero 09-06-2003}
\vskip1cm
\begin{enumerate}
\item Nota che
\[
\begin{split}
[e^{3n}+e^n]^{\frac 13}-e^n&=e^n\{[1+e^{-2n}]^{\frac 13}-1\}=
e^n\{\frac 13 e^{-2n}+\Or(e^{-4n})\}\\
&= \frac 13 e^{-n}+\Or (e^{-3n}).
\end{split}
\]
Dunque
\[
\lim_{n\to\infty}\left\{[e^{3n}+e^n]^{\frac 13}-e^n\right\}^{\frac
1n}=e^{-1}.
\]
Da cui segue che il raggio di convergenza \`e $e$.
\item Poich\`e per ogni $-1<a<1$ si ha
\[
\frac 1{1-a}=\sum_{n=0}^\infty a^n
\]
segue che, per ogni $|x|\leq 2^{-\frac 12}$,
\[
f(x)=\sum_{n=0}^\infty (-2)^nx^{2n},
\]
che \`e chiaramente la serie di Taylor cercata.
\item L'equazione si risolve per separazione di variabili:
\[
(x-ax^2)^{-1}x'=1
\]
che, integrata rispetto a $t$, da
\[
x_a(t)=\frac{e^t}{1-a+a e^t}.
\]
Dunque
\[
\lim_{t\to\infty} x_a(t)=\frac 1 a.
\]
\item Poich\`e $f$ \`e dispari tutti i coefficienti del coseno sono
nulli.
Inol\-tre
\[
\frac 1\pi\int_{-\pi}^\pi f(x)\sin nx \;dx=\frac
1\pi\left\{-\int_{-\pi}^0\sin nx+\int_0^\pi\sin nx\right\}=\frac
2{n\pi} (1-(-1)^n).
\]
Quindi,
\[
f(x)=\sum_{n=1}^\infty\frac {2(1-(-1)^n)}{n\pi}\sin nx
\]
e, ponendo $x=\frac \pi 2$, si ottiene
\[
1=\sum_{n=1}^\infty\frac {2(1-(-1)^n)}{n\pi}\sin
\frac{n\pi}2=\sum_{n=0}^\infty\frac{4(-1)^n}{\pi(2n+1)} 
\]
da cui segue
\[
\sum_{n=0}^\infty\frac{(-1)^n}{2n+1}=\frac \pi 4.
\]
\end{enumerate}
\end{document}
